% answers.tex -- answers and worked solutions for booklet.tex (Module 7 MC booklet), same
% layout (quiz.sty). One practiceans per quizq, in the same order (72).
% Every number: checks/m7-mc-booklet/numbers.py
% Build: TEXINPUTS=.: xelatex answers.tex (twice)
\documentclass[12pt]{article}
\usepackage{quiz}
\hypersetup{pdftitle={Multiple Choice Booklet - Module 7 - The Nature of Light - Answers}}
% links between the two PDF files are not possible, so question headings are plain text here
\renewcommand\hyperlink[2]{#2}
\newcommand\anstopic[1]{\par\Needspace{8\baselineskip}\vspace{4mm}%
  \begin{center}{\headfont\bfseries\fontsize{14}{18}\selectfont #1\par}\end{center}\vspace{2pt}}

\begin{document}
\quiztitle{Year 12 Physics}{Multiple Choice Booklet - Module 7 - The Nature of Light: Answers}

\begin{center}
  \renewcommand\arraystretch{1.45}
  \begin{tabular}{|l|*{12}{>{\centering\arraybackslash}p{7mm}|}}
    \hline
    \textbf{Q} & 1 & 2 & 3 & 4 & 5 & 6 & 7 & 8 & 9 & 10 & 11 & 12\\ \hline
    \textbf{Ans} & A & D & C & D & C & B & B & D & C & A & A & A\\ \hline\hline
    \textbf{Q} & 13 & 14 & 15 & 16 & 17 & 18 & 19 & 20 & 21 & 22 & 23 & 24\\ \hline
    \textbf{Ans} & D & B & B & B & C & A & B & B & C & C & B & C\\ \hline\hline
    \textbf{Q} & 25 & 26 & 27 & 28 & 29 & 30 & 31 & 32 & 33 & 34 & 35 & 36\\ \hline
    \textbf{Ans} & C & D & A & C & A & D & B & D & A & D & A & B\\ \hline\hline
    \textbf{Q} & 37 & 38 & 39 & 40 & 41 & 42 & 43 & 44 & 45 & 46 & 47 & 48\\ \hline
    \textbf{Ans} & A & C & A & B & B & C & D & B & C & A & B & C\\ \hline\hline
    \textbf{Q} & 49 & 50 & 51 & 52 & 53 & 54 & 55 & 56 & 57 & 58 & 59 & 60\\ \hline
    \textbf{Ans} & D & D & D & C & A & D & A & D & C & C & B & A\\ \hline\hline
    \textbf{Q} & 61 & 62 & 63 & 64 & 65 & 66 & 67 & 68 & 69 & 70 & 71 & 72\\ \hline
    \textbf{Ans} & B & D & C & B & D & A & C & D & B & A & C & D\\ \hline
  \end{tabular}
\end{center}
Topic 1: Questions 1--18. \quad Topic 2: Questions 19--36. \quad Topic 3: Questions 37--54. \quad
Topic 4: Questions 55--72.

Useful relationships: $c = f\lambda$, \ $\lambda_\mathrm{max} = \dfrac{b}{T}$, \ $d\sin\theta = m\lambda$,
\ $I = I_\mathrm{max}\cos^2\theta$, \ $E = hf$, \ $K_\mathrm{max} = hf - \phi$, \ $p = \dfrac{h}{\lambda}$,
\ $t = \gamma t_0$, \ $L = \dfrac{L_0}{\gamma}$, \ $p = \gamma m v$, \ $E = mc^2$, where
$\gamma = \dfrac{1}{\sqrt{1 - v^2/c^2}}$.

\begin{practiceanswers}

% =====================================================================
\anstopic{Topic 1 - Electromagnetic Spectrum}

% ---- Q1
\begin{practiceans}
  \correct{A}
  Maxwell's equations predicted waves of changing electric and magnetic fields. Their speed came out
  as $\dfrac{1}{\sqrt{\mu_0\varepsilon_0}}$, using two constants measured in experiments on magnetism
  ($\mu_0$) and electricity ($\varepsilon_0$). The answer matched the measured speed of light, so he
  concluded that light is an electromagnetic wave.
  \begin{itemize}
    \item \xmark\ B: Fizeau's value was used to \emph{compare} with the prediction, not to make it.
    \item \xmark\ C and D: neither was used in Maxwell's calculation; these effects later led to the
          quantum model.
  \end{itemize}
\end{practiceans}

% ---- Q2
\Needspace{9\baselineskip}
\begin{practiceans}
  \correct{D}
  An EM wave is a transverse wave: $\vec{E}$ and $\vec{B}$ oscillate in phase, at \ang{90} to each other
  and at \ang{90} to the direction of travel.
  \begin{itemize}
    \item \xmark\ A: the fields are perpendicular, not parallel.
    \item \xmark\ B: EM waves need no medium; they travel through a vacuum.
    \item \xmark\ C: all EM waves travel at $c$ in a vacuum.
  \end{itemize}
\end{practiceans}

% ---- Q3
\Needspace{9\baselineskip}
\begin{practiceans}
  \correct{C}
  An accelerating charge produces EM waves. An oscillating charge is always accelerating, so it makes
  changing electric and magnetic fields that travel outwards as a wave.
  \begin{itemize}
    \item \xmark\ A: a charge at rest makes only a steady electric field.
    \item \xmark\ B and D: a charge moving at a constant velocity, or a steady current, does not
          radiate. Only accelerating charges produce EM waves.
  \end{itemize}
\end{practiceans}

% ---- Q4
\Needspace{9\baselineskip}
\begin{practiceans}
  \correct{D}
  $\mu_0\varepsilon_0 = (4\pi\times10^{-7})(\num{8.854e-12}) = \num{1.113e-17}$, so
  $\dfrac{1}{\sqrt{\mu_0\varepsilon_0}} = \dfrac{1}{\num{3.336e-9}} = \num{3.0e8}$ (in \si{m.s^{-1}}:
  the speed of light).
  \begin{itemize}
    \item \xmark\ A: this is $\mu_0\varepsilon_0$. \quad \xmark\ B: this is $\sqrt{\mu_0\varepsilon_0}$.
    \item \xmark\ C: this is $\dfrac{1}{\mu_0\varepsilon_0} = c^2$ (forgot the square root).
  \end{itemize}
\end{practiceans}

% ---- Q5
\Needspace{12\baselineskip}
\begin{practiceans}
  \correct{C}
  One gap to the next tooth is $\frac{1}{720}$ of a turn, so half of it is $\frac{1}{1440}$ of a turn.
  At 12.6 turns per second:
  \[ t = \frac{1}{1440 \times 12.6} = \qty{5.51e-5}{s}, \qquad
     c = \frac{2d}{t} = \frac{2 \times 8633}{\num{5.51e-5}} = \qty{3.1e8}{m.s^{-1}} \]
  \begin{itemize}
    \item \xmark\ A: forgot the 720 teeth (used the time for half a whole turn).
    \item \xmark\ B: used only one way ($d$, not $2d$). \quad \xmark\ D: counted the distance twice
          ($4d$).
  \end{itemize}
\end{practiceans}

% ---- Q6
\Needspace{9\baselineskip}
\begin{practiceans}
  \correct{B}
  A hot, low-pressure gas (hydrogen discharge tube) gives emission lines. A hot solid (filament) gives
  a continuous spectrum. Sunlight is a continuous spectrum with dark absorption lines, made by cooler
  gas in the Sun's outer layers.
  \begin{itemize}
    \item \xmark\ A, C and D: each has at least one source matched with the wrong type of spectrum.
  \end{itemize}
\end{practiceans}

% ---- Q7
\Needspace{9\baselineskip}
\begin{practiceans}
  \correct{B}
  The hot, dense interior emits a continuous spectrum. Atoms in the cooler outer layers absorb photons
  whose energies match their energy-level gaps, removing those wavelengths, so dark lines appear.
  \begin{itemize}
    \item \xmark\ A: the core is very hot; the lines come from absorption, not from missing emission.
    \item \xmark\ C: Earth's atmosphere does add some lines, but the star's own lines are made in the
          star.
    \item \xmark\ D: gravity does not slow light down in this way.
  \end{itemize}
\end{practiceans}

% ---- Q8
\Needspace{9\baselineskip}
\begin{practiceans}
  \correct{D}
  An element is present if \emph{all} of its lines appear in the star's spectrum. All the lines of X
  and all the lines of Z line up with dark lines in the star. None of the lines of Y do. So X and Z are
  present.
  \begin{itemize}
    \item \xmark\ A: the star also has two lines that match Z.
    \item \xmark\ B and C: no line of Y matches a line in the star.
  \end{itemize}
\end{practiceans}

% ---- Q9
\Needspace{9\baselineskip}
\begin{practiceans}
  \correct{C}
  Wien's law: $T = \dfrac{b}{\lambda_\mathrm{max}} = \dfrac{\num{2.898e-3}}{\num{290e-9}} = \qty{1.0e4}{K}$.
  \begin{itemize}
    \item \xmark\ A: this is $b \times \lambda_\mathrm{max}$.
    \item \xmark\ B: did not change nanometres into metres.
    \item \xmark\ D: wrong power of ten.
  \end{itemize}
\end{practiceans}

% ---- Q10
\Needspace{9\baselineskip}
\begin{practiceans}
  \correct{A}
  $\lambda_\mathrm{max} = \dfrac{b}{T}$: a shorter peak wavelength means a higher temperature. Star A
  peaks at the shorter wavelength (\qty{450}{nm}, in the blue part of the spectrum), so it is hotter
  and looks more blue. Star B peaks in the red, so it looks more red.
  \begin{itemize}
    \item \xmark\ B, C and D: each has at least one part the wrong way round.
  \end{itemize}
\end{practiceans}

% ---- Q11
\Needspace{9\baselineskip}
\begin{practiceans}
  \correct{A}
  All the lines are shifted to longer wavelengths (red shift), so the star is moving away from Earth
  (Doppler effect).
  \begin{itemize}
    \item \xmark\ B: moving towards Earth would give a blue shift (shorter wavelengths).
    \item \xmark\ C: temperature changes the peak of the continuous spectrum, not the positions of the
          lines.
    \item \xmark\ D: the same lines are seen, so the elements are the same; they are just shifted.
  \end{itemize}
\end{practiceans}

% ---- Q12
\Needspace{9\baselineskip}
\begin{practiceans}
  \correct{A}
  One side of a rotating star moves towards us and the other side moves away. The light from each part
  is Doppler-shifted by a different amount, so each line is spread over a range of wavelengths.
  \begin{itemize}
    \item \xmark\ B: rotation does not change the temperature.
    \item \xmark\ C: rotation does not make the gas denser.
    \item \xmark\ D: a shift of the whole spectrum comes from the motion of the whole star.
  \end{itemize}
\end{practiceans}

% ---- Q13
\Needspace{9\baselineskip}
\begin{practiceans}
  \correct{D}
  In a dense gas the atoms collide often, which broadens the lines. Narrow lines mean low-density gas
  in the outer layers. Giant stars have very low-density outer layers, so star P is likely a giant.
  \begin{itemize}
    \item \xmark\ A: motion away shifts the lines; it does not make them narrower.
    \item \xmark\ B: the stars have the same temperature.
    \item \xmark\ C: fewer elements would mean fewer lines, not narrower lines.
  \end{itemize}
\end{practiceans}

% ---- Q14
\Needspace{9\baselineskip}
\begin{practiceans}
  \correct{B}
  The pulse goes there and back, so the one-way time is $\frac{2.56}{2} = \qty{1.28}{s}$:
  $d = c \times 1.28 = (\num{3.00e8})(1.28) = \qty{3.84e8}{m}$.
  \begin{itemize}
    \item \xmark\ A: this is $\frac{c}{t}$. \quad \xmark\ C: forgot to halve the time (this is the
          round-trip distance).
    \item \xmark\ D: doubled instead of halved.
  \end{itemize}
\end{practiceans}

% ---- Q15
\Needspace{9\baselineskip}
\begin{practiceans}
  \correct{B}
  In order of increasing frequency (decreasing wavelength): radio, microwaves, infrared, visible,
  ultraviolet, X-rays, gamma rays.
  \begin{itemize}
    \item \xmark\ A: this is the order of \emph{decreasing} frequency.
    \item \xmark\ C and D: ultraviolet and visible light, or microwaves and visible light, are in the
          wrong order.
  \end{itemize}
\end{practiceans}

% ---- Q16
\Needspace{9\baselineskip}
\begin{practiceans}
  \correct{B}
  $\lambda = \dfrac{c}{f} = \dfrac{\num{3.00e8}}{\num{98.5e6}} = \qty{3.05}{m}$.
  \begin{itemize}
    \item \xmark\ A: this is $\frac{f}{c}$ (upside down). \quad \xmark\ C: did not change MHz into Hz.
    \item \xmark\ D: this is $c \times f$.
  \end{itemize}
\end{practiceans}

% ---- Q17
\Needspace{9\baselineskip}
\begin{practiceans}
  \correct{C}
  Hertz made radio waves with a spark gap and detected them with a receiver loop. He measured their
  speed (equal to $c$) and showed they reflect, refract, interfere and can be polarised, just like
  light. This confirmed Maxwell's prediction.
  \begin{itemize}
    \item \xmark\ A: the photon idea came from Einstein (Hertz did notice the photoelectric effect, but
          could not explain it).
    \item \xmark\ B: the charge on the electron was measured by Millikan.
    \item \xmark\ D: the speed of light does not depend on the observer's motion.
  \end{itemize}
\end{practiceans}

% ---- Q18
\Needspace{9\baselineskip}
\begin{practiceans}
  \correct{A}
  Each element has its own set of energy levels. Photons are emitted or absorbed only when electrons
  move between these levels, so each element has its own pattern of lines, like a fingerprint.
  \begin{itemize}
    \item \xmark\ B: gases (not solids) give line spectra; hot solids give continuous spectra.
    \item \xmark\ C: there is no such simple rule.
    \item \xmark\ D: different elements produce lines at different wavelengths.
  \end{itemize}
\end{practiceans}

% =====================================================================
\anstopic{Topic 2 - Light: Wave Model}

% ---- Q19
\begin{practiceans}
  \correct{B}
  The narrower the slit compared with the wavelength, the more the light spreads out. So the central
  bright band becomes wider.
  \begin{itemize}
    \item \xmark\ A and C: diffraction increases as the slit gets narrower.
    \item \xmark\ D: light still passes through and diffracts.
  \end{itemize}
\end{practiceans}

% ---- Q20
\Needspace{9\baselineskip}
\begin{practiceans}
  \correct{B}
  For small angles, the fringe spacing is
  $\Delta y = \dfrac{\lambda L}{d} = \dfrac{(\num{600e-9})(2.0)}{\num{0.30e-3}} = \qty{4.0e-3}{m} = \qty{4.0}{mm}$.
  \begin{itemize}
    \item \xmark\ A: half the spacing (centre to the first dark fringe).
    \item \xmark\ C: this is the distance from the centre to the second bright fringe.
    \item \xmark\ D: wrong power of ten (millimetres changed into metres wrongly).
  \end{itemize}
\end{practiceans}

% ---- Q21
\Needspace{9\baselineskip}
\begin{practiceans}
  \correct{C}
  $d = \dfrac{1}{600 \times 10^{3}} = \qty{1.67e-6}{m}$. \quad
  $\sin\theta = \dfrac{\lambda}{d} = \dfrac{\num{550e-9}}{\num{1.67e-6}} = 0.33$, so $\theta = \ang{19}$.
  \begin{itemize}
    \item \xmark\ A: this is the value of $\sin\theta$, not the angle.
    \item \xmark\ B: used half the correct value of $\frac{\lambda}{d}$.
    \item \xmark\ D: this is the second-order angle.
  \end{itemize}
\end{practiceans}

% ---- Q22
\Needspace{9\baselineskip}
\begin{practiceans}
  \correct{C}
  $\sin\theta$ cannot be more than 1, so $m \le \dfrac{d}{\lambda} = \dfrac{\num{1.67e-6}}{\num{550e-9}} = 3.03$.
  The highest whole-number order is 3.
  \begin{itemize}
    \item \xmark\ A and B: orders 2 and 3 also exist ($\sin\theta = 0.66$ and $0.99$).
    \item \xmark\ D: order 4 would need $\sin\theta = 1.32$, which is impossible.
  \end{itemize}
\end{practiceans}

% ---- Q23
\Needspace{9\baselineskip}
\begin{practiceans}
  \correct{B}
  $\Delta y = \dfrac{\lambda L}{d}$. Red light has a longer wavelength than green light, so the fringes
  are further apart.
  \begin{itemize}
    \item \xmark\ A: this would happen for a shorter wavelength (for example blue).
    \item \xmark\ C: the spacing depends on $\lambda$. \quad \xmark\ D: red light still interferes.
  \end{itemize}
\end{practiceans}

% ---- Q24
\Needspace{9\baselineskip}
\begin{practiceans}
  \correct{C}
  With one slit, there is no second source to interfere with, so the closely spaced two-slit fringes
  disappear. Only the single-slit diffraction pattern is left: a broad, bright central band with weak
  bands on each side.
  \begin{itemize}
    \item \xmark\ A: light still passes through the open slit.
    \item \xmark\ B: the two-slit fringes need light from both slits.
    \item \xmark\ D: the fringe spacing depends on the slit separation, which no longer exists.
  \end{itemize}
\end{practiceans}

% ---- Q25
\Needspace{9\baselineskip}
\begin{practiceans}
  \correct{C}
  Dark fringes are where the path difference is $(m + \frac{1}{2})\lambda$: $\frac{1}{2}\lambda$ for the
  first, $\frac{3}{2}\lambda$ for the second.
  \begin{itemize}
    \item \xmark\ A: this is the first dark fringe.
    \item \xmark\ B and D: whole numbers of wavelengths give bright fringes.
  \end{itemize}
\end{practiceans}

% ---- Q26
\Needspace{9\baselineskip}
\begin{practiceans}
  \correct{D}
  Huygens' wave model predicted that light slows down in a denser medium such as water. Newton's
  particle model needed the particles to speed up in water to explain refraction. Foucault measured a
  lower speed in water, so this supported the wave model.
  \begin{itemize}
    \item \xmark\ A and B: Newton's model predicted that light \emph{speeds up} in water.
    \item \xmark\ C: the wave model predicted a lower speed.
  \end{itemize}
\end{practiceans}

% ---- Q27
\Needspace{9\baselineskip}
\begin{practiceans}
  \correct{A}
  Newton said the particles of light are attracted by the denser medium. This increases the part of
  their velocity perpendicular to the surface, so they bend towards the normal and travel faster in
  glass.
  \begin{itemize}
    \item \xmark\ B: in Newton's model the particles speed up in glass and bend towards the normal.
    \item \xmark\ C: this is the wave (Huygens) explanation.
    \item \xmark\ D: Newton's particles did not interfere.
  \end{itemize}
\end{practiceans}

% ---- Q28
\Needspace{9\baselineskip}
\begin{practiceans}
  \correct{C}
  The first filter halves unpolarised light: $\frac{I_0}{2}$. Malus' law for the second filter:
  $I = \dfrac{I_0}{2}\cos^2 60^\circ = \dfrac{I_0}{2} \times \dfrac{1}{4} = \dfrac{I_0}{8}$.
  \begin{itemize}
    \item \xmark\ A: the intensity after the first filter only.
    \item \xmark\ B: forgot that the first filter halves the intensity.
    \item \xmark\ D: used $\sin^2 60^\circ$ instead of $\cos^2 60^\circ$.
  \end{itemize}
\end{practiceans}

% ---- Q29
\Needspace{9\baselineskip}
\begin{practiceans}
  \correct{A}
  $\dfrac{I}{I_\mathrm{max}} = \cos^2\theta = 0.75$, so $\cos\theta = 0.866$ and $\theta = \ang{30}$.
  \begin{itemize}
    \item \xmark\ B: $\cos^2 45^\circ = 0.50$ (50\%). \quad \xmark\ C: $\cos^2 60^\circ = 0.25$ (25\%).
    \item \xmark\ D: $\cos^2 75^\circ = 0.067$ (about 7\%).
  \end{itemize}
\end{practiceans}

% ---- Q30
\Needspace{9\baselineskip}
\begin{practiceans}
  \correct{D}
  Only transverse waves can be polarised, because their oscillations can be in different directions
  at \ang{90} to the direction of travel. Sound is longitudinal, so it cannot be polarised.
  \begin{itemize}
    \item \xmark\ A and C: these are true, but polarisation does not show them.
    \item \xmark\ B: polarisation is evidence for the wave model.
  \end{itemize}
\end{practiceans}

% ---- Q31
\Needspace{9\baselineskip}
\begin{practiceans}
  \correct{B}
  After filter 1: $\frac{I_0}{2}$. After filter 2 (at \ang{45}): $\frac{I_0}{2}\cos^2 45^\circ = \frac{I_0}{4}$.
  After filter 3 (another \ang{45}): $\frac{I_0}{4}\cos^2 45^\circ = \frac{I_0}{8}$.
  \begin{itemize}
    \item \xmark\ A: the middle filter turns the plane of polarisation, so some light gets through.
    \item \xmark\ C: used only one of the two $\cos^2 45^\circ$ factors.
    \item \xmark\ D: this is the intensity after the first filter only.
  \end{itemize}
\end{practiceans}

% ---- Q32
\Needspace{9\baselineskip}
\begin{practiceans}
  \correct{D}
  $I = I_\mathrm{max}\cos^2\theta$: maximum at \ang{0}, zero at \ang{90}, maximum again at \ang{180},
  with a smooth curved shape.
  \begin{itemize}
    \item \xmark\ A: the intensity is not a straight line, and it rises again after \ang{90}.
    \item \xmark\ B: this is $|\cos\theta|$ (it has a sharp point at \ang{90}).
    \item \xmark\ C: this is $\sin^2\theta$ (it starts at zero).
  \end{itemize}
\end{practiceans}

% ---- Q33
\Needspace{9\baselineskip}
\begin{practiceans}
  \correct{A}
  The maxima are still where $d\sin\theta = m\lambda$, so they are at the same angles. With many slits,
  the waves cancel almost everywhere else, so the maxima are much sharper, and more light gets through,
  so they are brighter.
  \begin{itemize}
    \item \xmark\ B: the same $d$ gives the same angles.
    \item \xmark\ C: more slits make the maxima narrower and brighter.
    \item \xmark\ D: higher orders still appear.
  \end{itemize}
\end{practiceans}

% ---- Q34
\Needspace{9\baselineskip}
\begin{practiceans}
  \correct{D}
  At the centre ($m = 0$), every wavelength is at $\theta = 0$, so the colours add up to white. In each
  spectrum, $\sin\theta = \frac{m\lambda}{d}$: violet (shortest $\lambda$) is closest to the centre and red is
  furthest.
  \begin{itemize}
    \item \xmark\ A and B: the central maximum is white. \quad \xmark\ C: the angle depends on
          $\lambda$, so the colours separate.
  \end{itemize}
\end{practiceans}

% ---- Q35
\Needspace{9\baselineskip}
\begin{practiceans}
  \correct{A}
  Diffraction is large when the wavelength is similar to (or larger than) the size of the gap or
  obstacle. Radio wavelengths are metres to hundreds of metres, so they bend around buildings. Light wavelengths are
  less than a millionth of a metre, so light hardly bends and the shadow is sharp.
  \begin{itemize}
    \item \xmark\ B: all EM waves travel at $c$.
    \item \xmark\ C: light does diffract, but only by a tiny amount around large objects.
    \item \xmark\ D: radio waves have \emph{lower} frequencies than light.
  \end{itemize}
\end{practiceans}

% ---- Q36
\Needspace{9\baselineskip}
\begin{practiceans}
  \correct{B}
  $\sin\theta = \dfrac{\lambda}{d} = \dfrac{\num{650e-9}}{\num{1.6e-6}} = 0.41$, so $\theta = \ang{24}$.
  \begin{itemize}
    \item \xmark\ A: this is the value of $\sin\theta$, not the angle.
    \item \xmark\ C: this is the second-order angle. \quad \xmark\ D: used $\cos^{-1}$ instead of $\sin^{-1}$.
  \end{itemize}
\end{practiceans}

% =====================================================================
\anstopic{Topic 3 - Light: Quantum Model}

% ---- Q37
\begin{practiceans}
  \correct{A}
  A hotter object emits more at every wavelength, so its whole curve is higher, and its peak is at a
  shorter wavelength ($\lambda_\mathrm{max} = \frac{b}{T}$). Only graph A shows both.
  \begin{itemize}
    \item \xmark\ B: the hotter curve peaks at a longer wavelength.
    \item \xmark\ C: the cooler curve is higher. \quad \xmark\ D: both curves peak at the same wavelength.
  \end{itemize}
\end{practiceans}

% ---- Q38
\Needspace{9\baselineskip}
\begin{practiceans}
  \correct{C}
  Classical physics (with energy shared equally among all possible waves) predicted that the intensity
  should keep rising as the wavelength gets shorter, giving infinite energy in the ultraviolet. Real
  black-body curves have a peak and fall to zero at short wavelengths. Planck solved this by
  quantising energy.
  \begin{itemize}
    \item \xmark\ A, B and D: none of these is what the name refers to.
  \end{itemize}
\end{practiceans}

% ---- Q39
\Needspace{9\baselineskip}
\begin{practiceans}
  \correct{A}
  $\lambda_\mathrm{max} = \dfrac{b}{T} = \dfrac{\num{2.898e-3}}{1100} = \qty{2.6e-6}{m}$ (infrared, with
  a little visible red light).
  \begin{itemize}
    \item \xmark\ B: wrong power of ten (Wien's constant in mm\,K). \quad \xmark\ C: this is $b \times T$.
    \item \xmark\ D: this is $\frac{T}{b}$.
  \end{itemize}
\end{practiceans}

% ---- Q40
\Needspace{9\baselineskip}
\begin{practiceans}
  \correct{B}
  $E = \dfrac{hc}{\lambda} = \dfrac{(\num{6.626e-34})(\num{3.00e8})}{\num{450e-9}} = \qty{4.4e-19}{J}$.
  \begin{itemize}
    \item \xmark\ A: this is $\frac{h}{\lambda}$ (the photon's momentum).
    \item \xmark\ C: wrong power of ten (nanometres changed into metres wrongly).
    \item \xmark\ D: this is the energy in electronvolts (\qty{2.8}{eV}), not joules.
  \end{itemize}
\end{practiceans}

% ---- Q41
\Needspace{9\baselineskip}
\begin{practiceans}
  \correct{B}
  Photon energy: $E = \dfrac{hc}{\lambda} = \dfrac{(\num{6.626e-34})(\num{3.00e8})}{\num{250e-9}}
  = \qty{7.95e-19}{J} = \qty{4.96}{eV}$. \quad $K_\mathrm{max} = E - \phi = 4.96 - 2.3 = \qty{2.7}{eV}$.
  \begin{itemize}
    \item \xmark\ A: this is the work function. \quad \xmark\ C: this is the photon energy.
    \item \xmark\ D: added the work function instead of subtracting it.
  \end{itemize}
\end{practiceans}

% ---- Q42
\Needspace{9\baselineskip}
\begin{practiceans}
  \correct{C}
  $\phi = 4.3 \times \num{1.602e-19} = \qty{6.89e-19}{J}$. \quad
  $f_0 = \dfrac{\phi}{h} = \dfrac{\num{6.89e-19}}{\num{6.626e-34}} = \qty{1.0e15}{Hz}$.
  \begin{itemize}
    \item \xmark\ A: this is the work function in joules. \quad \xmark\ B: this is the threshold
          \emph{wavelength} (in m).
    \item \xmark\ D: forgot to change eV into J.
  \end{itemize}
\end{practiceans}

% ---- Q43
\Needspace{9\baselineskip}
\begin{practiceans}
  \correct{D}
  $K_\mathrm{max} = hf - \phi$. The gradient is $h$ for every metal, so the two lines are parallel. Each
  metal has a different work function, so the lines cut the $f$ axis at different threshold frequencies,
  and there is no kinetic energy below the threshold.
  \begin{itemize}
    \item \xmark\ A and B: the gradients are different.
    \item \xmark\ C: no electrons are emitted at low frequencies, so the lines cannot start above zero
          at $f = 0$.
  \end{itemize}
\end{practiceans}

% ---- Q44
\Needspace{9\baselineskip}
\begin{practiceans}
  \correct{B}
  Brighter light of the same frequency means more photons per second, each with the same energy $hf$.
  More photons release more electrons per second, but each electron still gets the same energy, so
  $K_\mathrm{max} = hf - \phi$ does not change.
  \begin{itemize}
    \item \xmark\ A and C: $K_\mathrm{max}$ depends on frequency, not intensity.
    \item \xmark\ D: more photons means more electrons.
  \end{itemize}
\end{practiceans}

% ---- Q45
\Needspace{9\baselineskip}
\begin{practiceans}
  \correct{C}
  Each electron absorbs one photon. A red photon has less energy ($hf$) than the work function, so it
  cannot free an electron, no matter how many photons arrive. A violet photon has more energy than the
  work function.
  \begin{itemize}
    \item \xmark\ A: the bright red light has plenty of total energy, but it comes in photons that are
          each too small.
    \item \xmark\ B: this is not the reason. \quad \xmark\ D: waiting longer does not help; there is
          no build-up of energy.
  \end{itemize}
\end{practiceans}

% ---- Q46
\Needspace{9\baselineskip}
\begin{practiceans}
  \correct{A}
  $E = hf = (\num{6.626e-34})(\num{7.5e14}) = \qty{4.97e-19}{J} = \qty{3.1}{eV}$. \quad
  $K_\mathrm{max} = 3.1 - 2.0 = \qty{1.1}{eV}$. The stopping voltage gives $eV_0 = K_\mathrm{max}$, so
  $V_0 = \qty{1.1}{V}$.
  \begin{itemize}
    \item \xmark\ B: this is the work function. \quad \xmark\ C: this is the photon energy.
    \item \xmark\ D: added the work function instead of subtracting it.
  \end{itemize}
\end{practiceans}

% ---- Q47
\Needspace{9\baselineskip}
\begin{practiceans}
  \correct{B}
  $p = \dfrac{h}{\lambda} = \dfrac{\num{6.626e-34}}{\num{500e-9}} = \qty{1.3e-27}{kg.m.s^{-1}}$.
  \begin{itemize}
    \item \xmark\ A: this is $h \times \lambda$. \quad \xmark\ C: this is the photon's energy (in J).
    \item \xmark\ D: this is the frequency (in Hz).
  \end{itemize}
\end{practiceans}

% ---- Q48
\Needspace{9\baselineskip}
\begin{practiceans}
  \correct{C}
  Einstein said light is made of photons, each of energy $E = hf$. One photon gives all its energy to one
  electron. If $hf > \phi$, the electron escapes with $K_\mathrm{max} = hf - \phi$.
  \begin{itemize}
    \item \xmark\ A: this is the wave-model idea that failed.
    \item \xmark\ B and D: emission and $K_\mathrm{max}$ depend on frequency, not brightness.
  \end{itemize}
\end{practiceans}

% ---- Q49
\Needspace{9\baselineskip}
\begin{practiceans}
  \correct{D}
  In the wave model, energy is spread over the wave, so with dim light an electron would need time to
  build up enough energy. Electrons are actually emitted straight away, because a single photon gives
  all its energy at once.
  \begin{itemize}
    \item \xmark\ A: the wave model also predicts more electrons for brighter light.
    \item \xmark\ B and C: these do not contradict the wave model.
  \end{itemize}
\end{practiceans}

% ---- Q50
\Needspace{9\baselineskip}
\begin{practiceans}
  \correct{D}
  Both curves reach zero at the same stopping voltage $-V_0$, so the electrons have the same
  $K_\mathrm{max}$, meaning the frequency is the same. X has a larger maximum (saturation) current, so
  more electrons are released per second: X has a greater intensity.
  \begin{itemize}
    \item \xmark\ A and C: a higher frequency would need a larger stopping voltage.
    \item \xmark\ B: the same intensity would give the same saturation current.
  \end{itemize}
\end{practiceans}

% ---- Q51
\Needspace{9\baselineskip}
\begin{practiceans}
  \correct{D}
  Planck assumed that energy is quantised: the oscillators can only have energies $nhf$
  ($n = 1, 2, 3, \dots$). This made his formula match the measured black-body curves.
  \begin{itemize}
    \item \xmark\ A: this is Einstein's second postulate of special relativity.
    \item \xmark\ B: this is Bohr's model of the atom. \quad \xmark\ C: this is the classical wave idea.
  \end{itemize}
\end{practiceans}

% ---- Q52
\Needspace{9\baselineskip}
\begin{practiceans}
  \correct{C}
  Energy of one photon: $E = \dfrac{hc}{\lambda} = \dfrac{(\num{6.626e-34})(\num{3.00e8})}{\num{650e-9}}
  = \qty{3.06e-19}{J}$. \quad Number per second $= \dfrac{P}{E} = \dfrac{\num{5.0e-3}}{\num{3.06e-19}}
  = \num{1.6e16}$.
  \begin{itemize}
    \item \xmark\ A: this is the energy of one photon.
    \item \xmark\ B: used $\qty{5.0e-6}{W}$ instead of \qty{5.0e-3}{W}.
    \item \xmark\ D: used \qty{5.0}{W} (forgot the milli).
  \end{itemize}
\end{practiceans}

% ---- Q53
\Needspace{9\baselineskip}
\begin{practiceans}
  \correct{A}
  Work function: $\phi = \dfrac{hc}{\lambda_0} = \dfrac{(\num{6.626e-34})(\num{3.00e8})}{\num{500e-9}}
  = \qty{3.98e-19}{J} = \qty{2.5}{eV}$. \quad $K_\mathrm{max} = 3.5 - 2.5 = \qty{1.0}{eV}$.
  \begin{itemize}
    \item \xmark\ B: this is the work function. \quad \xmark\ C: this is the photon energy.
    \item \xmark\ D: added instead of subtracting.
  \end{itemize}
\end{practiceans}

% ---- Q54
\Needspace{9\baselineskip}
\begin{practiceans}
  \correct{D}
  $K_\mathrm{max} = hf - \phi$ has the form $y = mx + c$, so the gradient is $h$. The line cuts the $f$
  axis where $K_\mathrm{max} = 0$, which is the threshold frequency $f_0$. (The work function is the size
  of the intercept on the $K_\mathrm{max}$ axis, if the line is extended back.)
  \begin{itemize}
    \item \xmark\ A and C: the gradient is $h$, not $\phi$.
    \item \xmark\ B: the $f$-axis intercept is the threshold frequency, not the work function.
  \end{itemize}
\end{practiceans}

% =====================================================================
\anstopic{Topic 4 - Light and Special Relativity}

% ---- Q55
\begin{practiceans}
  \correct{A}
  Einstein's two postulates: (1) the laws of physics are the same in all inertial frames of reference;
  (2) the speed of light in a vacuum is the same for all observers in inertial frames.
  \begin{itemize}
    \item \xmark\ B and C: these are \emph{results} of special relativity, not postulates.
    \item \xmark\ D: special relativity removed the need for a medium (the aether).
  \end{itemize}
\end{practiceans}

% ---- Q56
\Needspace{9\baselineskip}
\begin{practiceans}
  \correct{D}
  Michelson and Morley expected the interference pattern to shift as their apparatus turned, because of
  Earth's motion through the aether. They found no shift (a null result), so no aether wind was
  detected.
  \begin{itemize}
    \item \xmark\ A and C: no difference in speed was detected.
    \item \xmark\ B: the experiment used the interference of light waves, and said nothing about
          particles.
  \end{itemize}
\end{practiceans}

% ---- Q57
\Needspace{9\baselineskip}
\begin{practiceans}
  \correct{C}
  The proper time is \qty{2.0}{s} (measured on the ship). $\gamma = \dfrac{1}{\sqrt{1 - 0.80^2}}
  = \dfrac{1}{0.60} = 1.67$, so $t = \gamma t_0 = 1.67 \times 2.0 = \qty{3.3}{s}$.
  \begin{itemize}
    \item \xmark\ A: divided by $\gamma$ instead of multiplying. \quad \xmark\ B: ignored time dilation.
    \item \xmark\ D: multiplied by $\gamma^2$.
  \end{itemize}
\end{practiceans}

% ---- Q58
\Needspace{9\baselineskip}
\begin{practiceans}
  \correct{C}
  $\gamma = \dfrac{1}{\sqrt{1 - 0.995^2}} = 10.0$. In Earth's frame the lifetime is
  $\gamma t_0 = 10.0 \times \qty{2.2}{\micro s} = \qty{22}{\micro s}$, so
  $d = vt = (0.995 \times \num{3.00e8})(\num{22e-6}) = \qty{6.6e3}{m} = \qty{6.6}{km}$.
  \begin{itemize}
    \item \xmark\ A: divided by $\gamma$ instead of multiplying. \quad \xmark\ B: ignored time dilation.
    \item \xmark\ D: applied $\gamma$ twice.
  \end{itemize}
\end{practiceans}

% ---- Q59
\Needspace{9\baselineskip}
\begin{practiceans}
  \correct{B}
  $\gamma = \dfrac{1}{\sqrt{1 - 0.60^2}} = \dfrac{1}{0.80} = 1.25$, so
  $L = \dfrac{L_0}{\gamma} = \dfrac{120}{1.25} = \qty{96}{m}$.
  \begin{itemize}
    \item \xmark\ A: multiplied by $0.60$ (the speed) instead of $0.80$. \quad \xmark\ C: this is the
          proper length.
    \item \xmark\ D: multiplied by $\gamma$ instead of dividing.
  \end{itemize}
\end{practiceans}

% ---- Q60
\Needspace{9\baselineskip}
\begin{practiceans}
  \correct{A}
  Proper time is measured by the observer for whom both events happen at the same place. The heartbeats
  happen at the same place in the astronaut's frame, so she measures the proper time. The Earth observer
  measures a longer (dilated) time.
  \begin{itemize}
    \item \xmark\ B: in the Earth frame the heartbeats happen at different places. Also, neither frame is
          ``really'' at rest.
    \item \xmark\ C: the two observers measure different times. \quad \xmark\ D: the astronaut does.
  \end{itemize}
\end{practiceans}

% ---- Q61
\Needspace{9\baselineskip}
\begin{practiceans}
  \correct{B}
  $\gamma = \dfrac{1}{\sqrt{1 - 0.90^2}} = 2.29$. \quad
  $p = \gamma m v = 2.29 \times (\num{9.109e-31}) \times (0.90 \times \num{3.00e8}) = \qty{5.6e-22}{kg.m.s^{-1}}$.
  \begin{itemize}
    \item \xmark\ A: this is the classical momentum $mv$ (no $\gamma$).
    \item \xmark\ C: multiplied by $\gamma^2$. \quad \xmark\ D: this is the speed $v$ (no mass).
  \end{itemize}
\end{practiceans}

% ---- Q62
\Needspace{9\baselineskip}
\begin{practiceans}
  \correct{D}
  $p = \gamma m v$. At low speeds $\gamma \approx 1$, so $p \approx mv$ (a straight line from the origin).
  As $v$ approaches $c$, $\gamma$ increases without limit, so the momentum rises steeply and has no
  maximum value. The speed never reaches $c$.
  \begin{itemize}
    \item \xmark\ A: this is classical momentum, and it goes past $c$.
    \item \xmark\ B: the momentum does not level off. \quad \xmark\ C: this shape goes past $c$ and
          is not linear at low speeds.
  \end{itemize}
\end{practiceans}

% ---- Q63
\Needspace{9\baselineskip}
\begin{practiceans}
  \correct{C}
  $E = mc^2 = (\num{3.0e-29})(\num{3.00e8})^2 = \qty{2.7e-12}{J}$.
  \begin{itemize}
    \item \xmark\ A: divided by $c^2$. \quad \xmark\ B: multiplied by $c$ instead of $c^2$.
    \item \xmark\ D: multiplied by $c^3$.
  \end{itemize}
\end{practiceans}

% ---- Q64
\Needspace{9\baselineskip}
\begin{practiceans}
  \correct{B}
  In one second the Sun radiates \qty{3.8e26}{J}. $m = \dfrac{E}{c^2} = \dfrac{\num{3.8e26}}{(\num{3.00e8})^2}
  = \qty{4.2e9}{kg}$ (about 4 million tonnes every second).
  \begin{itemize}
    \item \xmark\ A: this is $\frac{c^2}{E}$ (upside down). \quad \xmark\ C: divided by $c$ instead of
          $c^2$.
    \item \xmark\ D: multiplied by $c^2$ instead of dividing.
  \end{itemize}
\end{practiceans}

% ---- Q65
\Needspace{9\baselineskip}
\begin{practiceans}
  \correct{D}
  The satellites move fast relative to Earth, so by special relativity their clocks run slow (time
  dilation, about \qty{7}{\micro s} per day). GPS systems must correct for this, or position errors would
  build up quickly. (General relativity also makes them run fast because of weaker gravity, but the
  question asks about special relativity only.)
  \begin{itemize}
    \item \xmark\ A: the difference is small but important. \quad \xmark\ B: moving clocks run slow,
          not fast.
    \item \xmark\ C: the clocks never stop.
  \end{itemize}
\end{practiceans}

% ---- Q66
\Needspace{9\baselineskip}
\begin{practiceans}
  \correct{A}
  The flown clocks ended up different from the ground clocks by tiny amounts (nanoseconds), and the
  differences agreed with the predictions of special and general relativity. This was direct evidence
  of time dilation.
  \begin{itemize}
    \item \xmark\ B: the clocks did not stay in step.
    \item \xmark\ C: the eastward and westward flights gave different results (one lost time, the other
          gained time).
    \item \xmark\ D: the experiment measured time, and the speed of light is the same for all observers.
  \end{itemize}
\end{practiceans}

% ---- Q67
\Needspace{9\baselineskip}
\begin{practiceans}
  \correct{C}
  The passenger moves towards the light from the front strike and away from the light from the back
  strike, so the light from the front reaches them first. Light travels at $c$ in their frame too and
  they are halfway between the ends, so they conclude that the front strike happened first.
  Simultaneity is relative.
  \begin{itemize}
    \item \xmark\ A: events that are simultaneous in one frame are not simultaneous in another frame
          moving relative to it.
    \item \xmark\ B: the back strike is seen second. \quad \xmark\ D: the passenger sees both.
  \end{itemize}
\end{practiceans}

% ---- Q68
\Needspace{9\baselineskip}
\begin{practiceans}
  \correct{D}
  By the first postulate, the laws of physics are the same in all inertial frames, so no experiment
  inside a closed laboratory moving at constant velocity can tell that it is moving.
  \begin{itemize}
    \item \xmark\ A: the pendulum behaves normally in its own frame.
    \item \xmark\ B: the speed of light is $c$ in every direction in every inertial frame.
    \item \xmark\ C: the scientist sees their own clocks run normally.
  \end{itemize}
\end{practiceans}

% ---- Q69
\Needspace{9\baselineskip}
\begin{practiceans}
  \correct{B}
  By the second postulate, the speed of light is the same for all observers, whatever the motion of the
  source. Light speeds do not add like ordinary speeds.
  \begin{itemize}
    \item \xmark\ A: subtracted the speeds. \quad \xmark\ C: added the speeds (Galilean addition).
    \item \xmark\ D: this has no physical basis.
  \end{itemize}
\end{practiceans}

% ---- Q70
\Needspace{9\baselineskip}
\begin{practiceans}
  \correct{A}
  Length contraction happens only along the direction of motion. The length in that direction is
  divided by $\gamma$ ($\gamma = 2.3$ at $0.90c$); the height and width do not change.
  \begin{itemize}
    \item \xmark\ B: perpendicular lengths are not contracted. \quad \xmark\ C: moving objects are
          shorter, not longer.
    \item \xmark\ D: at $0.90c$ the contraction is large.
  \end{itemize}
\end{practiceans}

% ---- Q71
\Needspace{9\baselineskip}
\begin{practiceans}
  \correct{C}
  Half the rate means $\gamma = 2$. $\sqrt{1 - \frac{v^2}{c^2}} = \frac{1}{2}$, so $\frac{v^2}{c^2} = \frac{3}{4}$
  and $v = 0.87c$.
  \begin{itemize}
    \item \xmark\ A: at $0.50c$, $\gamma = 1.15$. \quad \xmark\ B: at $0.71c$, $\gamma = 1.42$.
    \item \xmark\ D: at $0.97c$, $\gamma = 4.1$.
  \end{itemize}
\end{practiceans}

% ---- Q72
\Needspace{9\baselineskip}
\begin{practiceans}
  \correct{D}
  As the energy increases, $\gamma$ increases. The speed approaches $c$ but never reaches it, because
  that would need infinite energy. The energy ($\gamma mc^2$) and momentum ($\gamma mv$) can keep
  increasing without limit.
  \begin{itemize}
    \item \xmark\ A and C: nothing with mass can reach or go past $c$.
    \item \xmark\ B: the momentum keeps increasing as $\gamma$ increases.
  \end{itemize}
\end{practiceans}

\end{practiceanswers}
\end{document}
